% Vietnamese translation for OI Public Library
% Source-PDF: IOI中国国家候选队论文/2009/武森/浅谈信息学竞赛中的“0”和“1”.pdf
\documentclass[11pt]{article}
\usepackage{oipl-vn}

\newcommand{\Oh}{\mathcal{O}}
\newcommand{\lowbit}{\operatorname{lowbit}}
\newcommand{\xor}{\mathbin{\mathrm{xor}}}

\title{Bàn về ``0'' và ``1'' trong thi Tin học}
\author{Wu Sen\\Trường Trung học số 2 Thạch Gia Trang, Hà Bắc}
\date{2009}

\begin{document}
\maketitle

\origtitle{Bàn về ``0'' và ``1'' trong thi Tin học}
\sourcepath{IOI China candidate team papers / 2009 / Wu Sen}

\begin{abstract}
Máy tính hiện đại biểu diễn và xử lý thông tin bằng hệ nhị phân. Không chỉ
trong kiến trúc máy tính, tư tưởng nhị phân còn xuất hiện rất nhiều trong
thi Tin học: từ cấu trúc dữ liệu, nén trạng thái, mô hình hóa bằng bit cho
tới cây nhị phân \(0/1\). Bài viết dùng bốn ví dụ để tổng kết cách đưa hai
ký hiệu \(0\) và \(1\) vào phân tích thuật toán.
\end{abstract}

\paragraph{Từ khóa.}
Nhị phân; thập phân; cây Fenwick; nén trạng thái; cây nhị phân \(0/1\).

\tableofcontents

\section{Ứng dụng trong cấu trúc dữ liệu}

\subsection{Ví dụ 1: Matrix}

Cho ma trận \(M\times N\), mỗi ô chỉ có giá trị \(0\) hoặc \(1\). Ban đầu mọi
ô đều bằng \(0\). Mỗi thao tác cập nhật chọn một hình chữ nhật có góc trái
trên \((x_1,y_1)\) và góc phải dưới \((x_2,y_2)\), rồi đảo bit mọi ô trong
hình chữ nhật. Mỗi truy vấn hỏi giá trị ô \((x,y)\).

Vì mỗi ô chỉ là \(0\) hoặc \(1\), ta chỉ cần biết ô đó bị đảo bao nhiêu lần;
đáp án là số lần đảo modulo \(2\).

\subsubsection{Từ một chiều tới hai chiều}

Trước hết xét mảng một chiều \(d\) độ dài \(N\). Một cập nhật đoạn \([x,y]\)
có thể được đổi thành hai điểm mốc:
\[
  d_x \mathrel{+}=1,\qquad d_{y+1}\mathrel{+}=1.
\]
Khi hỏi vị trí \(x\), tính
\[
  \operatorname{Sum}_x=\sum_{i=1}^{x} d_i,
\]
rồi trả \(\operatorname{Sum}_x\bmod 2\).

Với hai chiều, một hình chữ nhật được đánh dấu bằng bốn góc:
\[
  (x_1,y_1),\quad (x_2+1,y_1),\quad
  (x_1,y_2+1),\quad (x_2+1,y_2+1).
\]
Khi truy vấn \((x,y)\), tính tổng tiền tố hai chiều
\[
  \operatorname{Sum}_{x,y}
  =\sum_{i=1}^{x}\sum_{j=1}^{y} d_{i,j},
\]
và lấy modulo \(2\). Bốn điểm mốc làm cho các vùng ngoài hình chữ nhật nhận
thêm \(0\), \(2\) hoặc \(4\) lần, nên không đổi modulo \(2\); các điểm bên
trong nhận thêm đúng \(1\) lần. Cách chứng minh tương tự như hiệu phân nhiều
chiều.

Lập luận này mở rộng tự nhiên: trong ba chiều cần \(8\) điểm mốc, còn trong
\(n\) chiều cần \(2^n\) điểm quanh khối cần cập nhật.

\subsubsection{Cây Fenwick và tư tưởng nhị phân}

Để duy trì tổng tiền tố, một cấu trúc ngắn gọn và dễ mở rộng là cây Fenwick
(Binary Indexed Tree). Mỗi chỉ số được nhìn dưới dạng nhị phân. Nếu
\[
  \lowbit(x)=x\mathbin{\&}(-x),
\]
thì nút \(x\) lưu tổng của đoạn độ dài \(\lowbit(x)\) kết thúc tại \(x\).
Ví dụ, \(11=(1011)_2\), và phần bit thấp nhất quyết định kích thước đoạn mà
nút này đại diện.

\paragraph{Cập nhật.}
\begin{verbatim}
while x <= max:
    C[x] += delta
    x += x & -x
\end{verbatim}

\paragraph{Truy vấn tiền tố.}
\begin{verbatim}
sum = 0
while x > 0:
    sum += C[x]
    x -= x & -x
\end{verbatim}

Trong truy vấn, thao tác \(x-\lowbit(x)\) xóa bit \(1\) thấp nhất của \(x\).
Một số không vượt quá \(n\) có nhiều nhất \(\log n\) bit \(1\), nên truy vấn
mất \(\Oh(\log n)\). Cập nhật đi qua các tổ tiên tương ứng và cũng mất
\(\Oh(\log n)\). Ưu điểm của Fenwick là mã ngắn, khó sai, có thể mở rộng lên
hai hoặc ba chiều, và rất phù hợp với các bài mà thao tác không quá phức tạp.

\section{Ứng dụng trong tư duy giải bài}

\subsection{Ví dụ 2: Sudoku}

Sudoku là bảng \(9\times 9\), chia thành chín khối \(3\times 3\). Cần điền
các số \(1\) đến \(9\) sao cho mỗi hàng, mỗi cột và mỗi khối đều chứa đủ
chín số đúng một lần.

Cách tìm kiếm trực tiếp theo thứ tự ô thường quá chậm. Một tối ưu quen thuộc
là với mỗi ô trống, ghi lại tập chữ số còn có thể đặt, chọn ô có ít lựa chọn
nhất để rẽ nhánh, rồi cập nhật các ô cùng hàng, cột và khối. Tuy nhiên nếu
cài bằng tập hợp thông thường, thao tác sửa và sắp lại ứng viên khá rườm rà.

Tư tưởng nhị phân biến tập ứng viên của một ô thành một số \(9\) bit. Bit thứ
\(i\) biểu thị chữ số \(i\) có bị cấm hay không; chẳng hạn \(0\) nghĩa là có
thể đặt, \(1\) nghĩa là không thể đặt. Mỗi ô khi đó được biểu diễn bởi một
số từ \(0\) tới \(511\). Trạng thái hàng, cột và khối cũng có thể là mặt nạ
bit:
\[
  \text{available}(r,c)=\neg(row_r\lor col_c\lor block_{r,c}) \mathbin{\&} 511.
\]
Khi tìm kiếm, ta duyệt các bit \(1\) còn đặt được bằng các phép bit. Chỉ riêng
nén trạng thái bằng bit đã đủ làm bài toán nhanh hơn nhiều; kết hợp thêm chọn
ô có ít ứng viên nhất sẽ cho lời giải rất ổn định.

\subsection{Ví dụ 3: Requirements}

Cho \(N\) điểm năm chiều
\[
  A=(x_1,x_2,x_3,x_4,x_5),
\]
cần tìm hai điểm \(X,Y\) có khoảng cách Manhattan lớn nhất:
\[
  \sum_{i=1}^{5}|x_i-y_i|.
\]
Vì \(N\) có thể tới \(100000\), không thể xét mọi cặp điểm.

Trước hết xét hai chiều. Với hai điểm, đặt \(a=x_i-y_i\), \(b=x_j-y_j\). Ta
có
\[
|a|+|b|
=\max\{a+b,\ a-b,\ -a+b,\ -a-b\}.
\]
Nói cách khác, ta chỉ cần liệt kê dấu cộng/trừ trước từng tọa độ. Mỗi cách
chọn dấu có thể mã hóa bằng một mặt nạ nhị phân: bit \(1\) nghĩa là dấu cộng,
bit \(0\) nghĩa là dấu trừ. Với mỗi mặt nạ \(m\), tính giá trị tuyến tính
\[
  F_m(P)=\sum_{i=1}^{d} s_i x_i,\qquad s_i\in\{-1,1\},
\]
rồi lưu \(\max F_m(P)\) và \(\min F_m(P)\) trên mọi điểm. Đáp án là
\[
  \max_m \bigl(\max_P F_m(P)-\min_P F_m(P)\bigr).
\]
Với \(d=5\), chỉ có \(2^5=32\) mặt nạ, nên thuật toán chạy trong
\(\Oh(2^dN)\), rất nhỏ so với \(\Oh(N^2)\). Điều kiện quan trọng là số chiều
nhỏ hơn nhiều so với số điểm.

\subsection{Ví dụ 4: Cow XOR}

Cho dãy \(a_1,a_2,\ldots,a_N\), cần chọn một đoạn liên tiếp sao cho xor của
đoạn là lớn nhất. Nếu có nhiều đoạn, chọn đoạn có chỉ số kết thúc nhỏ nhất;
nếu vẫn hòa, chọn đoạn ngắn nhất. Đặt
\[
  S_k=a_1\xor a_2\xor\cdots\xor a_k.
\]
Khi đó
\[
  a_i\xor a_{i+1}\xor\cdots\xor a_j=S_j\xor S_{i-1}.
\]
Vì các số nằm trong khoảng \(0\) tới \(2^{21}-1\), ta có thể lưu các tiền tố
\(S_i\) trong cây nhị phân \(0/1\) độ cao \(21\). Mỗi cạnh ứng với một bit
của số; lá lưu chỉ số tiền tố.

Khi đang xét \(S_i\), để tìm tiền tố \(S_t\) làm \(S_i\xor S_t\) lớn nhất,
ta đi từ bit cao xuống thấp. Nếu bit hiện tại của \(S_i\) là \(0\), ưu tiên
nhánh \(1\); nếu là \(1\), ưu tiên nhánh \(0\). Tức là luôn cố làm bit cao
nhất còn lại của xor bằng \(1\). Chiến lược tham lam đúng vì một bit cao
\(2^k\) lớn hơn tổng mọi bit thấp hơn:
\[
  2^k>\sum_{i=0}^{k-1}2^i.
\]
Mỗi lần truy vấn và chèn một tiền tố mất \(\Oh(21)\), nên toàn bộ thuật toán
là \(\Oh(N\log A)\), với \(A<2^{21}\).

\section{Tổng kết}

Các ví dụ trên cho thấy tư tưởng nhị phân không chỉ là cách biểu diễn số.
Trong cấu trúc dữ liệu, nó tạo ra Fenwick tree với mã ngắn và hiệu quả tốt.
Trong giải bài, nó giúp nén trạng thái, liệt kê tổ hợp dấu, xây mô hình bằng
mặt nạ bit và dùng cây nhị phân \(0/1\) để tối ưu xor. Khi biết chuyển bài
toán từ hệ thập phân quen thuộc sang góc nhìn nhị phân, nhiều lời giải trở
nên ngắn hơn, rõ hơn và nhanh hơn.

\section{Tài liệu tham khảo}

\begin{enumerate}
  \item Liu Rujia và Huang Liang, \textit{Nghệ thuật thuật toán và thi Tin
  học}, Nhà xuất bản Đại học Thanh Hoa, ấn bản tháng 1 năm 2004.
  \item Peking University Judge Online, \texttt{http://acm.pku.edu.cn/}.
  \item Zhejiang University Judge Online, \texttt{http://acm.zju.edu.cn/}.
  \item USA Computing Olympiad, \texttt{http://ace.delos.com/}.
\end{enumerate}

\section{Phụ lục: nguyên đề các ví dụ}

\subsection{Matrix}

Cho ma trận \(N\times N\) gồm các phần tử \(0\) hoặc \(1\), ban đầu đều là
\(0\). Có hai loại lệnh:
\begin{itemize}
  \item \texttt{C x1 y1 x2 y2}: đảo bit mọi ô trong hình chữ nhật có góc trái
  trên \((x_1,y_1)\) và góc phải dưới \((x_2,y_2)\).
  \item \texttt{Q x y}: hỏi giá trị \(A[x,y]\).
\end{itemize}
Có nhiều bộ test, \(N\le 1000\), số lệnh \(T\le 50000\).

\paragraph{Ví dụ.}
\begin{verbatim}
Input
1

2 10
C 2 1 2 2
Q 2 2
C 2 1 2 1
Q 1 1
C 1 1 2 1
C 1 2 1 2
C 1 1 2 2
Q 1 1
C 1 1 2 1
Q 2 1

Output
1
0
0
1
\end{verbatim}

\subsection{Sudoku}

Mỗi test là một dòng \(81\) ký tự mô tả bảng Sudoku theo thứ tự từng hàng.
Ký tự là chữ số \(1\) đến \(9\) hoặc dấu chấm cho ô trống. Mỗi bảng có đúng
một nghiệm; input kết thúc bằng dòng \texttt{end}. Với mỗi test, in bảng đã
hoàn thành trên một dòng.

\subsection{Requirements}

Mỗi trường đại học được mô tả bằng năm số thực không âm, tương ứng năm yêu
cầu tốt nghiệp. Độ khác nhau giữa hai trường \(X,Y\) là
\[
  |x_1-y_1|+|x_2-y_2|+|x_3-y_3|+|x_4-y_4|+|x_5-y_5|.
\]
Với \(1\le N\le 100000\), hãy in độ khác nhau lớn nhất giữa hai trường, làm
tròn đúng hai chữ số thập phân. Input kết thúc bằng \(N=0\).

\paragraph{Ví dụ.}
\begin{verbatim}
Input
3
2 5 6 2 1.5
1.2 3 2 5 4
7 5 3 2 5
0

Output
12.80
\end{verbatim}

\subsection{Cow XOR}

Cho \(N\) con bò theo thứ tự thứ bậc xã hội, mỗi con có một số trong khoảng
\(0\) tới \(2^{21}-1\). Chọn một đoạn liên tiếp sao cho xor các số trên đoạn
là lớn nhất. Nếu có nhiều đoạn, chọn đoạn có điểm kết thúc nhỏ nhất; nếu vẫn
hòa, chọn đoạn ngắn nhất. In giá trị xor lớn nhất, vị trí bắt đầu và vị trí
kết thúc.

\paragraph{Ví dụ.}
\begin{verbatim}
Input
5
1
0
5
4
2

Output
6 4 5
\end{verbatim}

\end{document}
